% Generated draft: requires mathematical and source review.
\documentclass[11pt]{article}
\usepackage[margin=25mm]{geometry}
\usepackage{fontspec}
\setmainfont{Latin Modern Roman}
\usepackage{amsmath,amssymb,mathtools,unicode-math}
\setmathfont{Latin Modern Math}
\usepackage{enumitem,needspace,xurl,hyperref}
\setlength{\parindent}{0pt}
\setlength{\parskip}{5pt}
\setlength{\emergencystretch}{3em}
\newcommand{\N}{\mathbb{N}}
\newcommand{\Z}{\mathbb{Z}}
\newcommand{\Q}{\mathbb{Q}}
\newcommand{\R}{\mathbb{R}}
\newcommand{\C}{\mathbb{C}}
\newcommand{\F}{\mathbb{F}}
\newcommand{\A}{\mathbb{A}}
\newcommand{\PP}{\mathbb{P}}
\newcommand{\E}{\mathbb{E}}
\newcommand{\eps}{\varepsilon}
\newcommand{\dd}{\,\mathrm{d}}
\begin{document}
\section*{UnsolvedMath ID: 30000891}
\textbf{Source problem code:} catalogue entry
\subsection{Definitions and mathematical statement}
Let $A\subset\N^m$ be a finite exponent set of cardinality $n$, contained in an affine hyperplane, and let $f(x)=(x^\alpha)_{\alpha\in A}$, where $x^\alpha=\prod_{r=1}^m x_r^{\alpha_r}$. The original source formulation takes $C$ to be the affine cone over the corresponding monomial image. With this cone convention, its $(k-1)$-st secant cone, for an integer $k\geq1$, is
\[
kC=\overline{\{c_1+\cdots+c_k:c_i\in C\}}^{\,\mathrm{Zar}}.
\]
One may construct $C$ as the cone over the Zariski closure of the projective monomial image of the algebraic torus. The ground field and characteristic are not supplied and require confirmation.

For real affine-linear forms $\ell_i:\R^m\to\R$, $1\leq i\leq k$, define
\[
C_i(\ell)=\{\alpha\in A:\ell_i(\alpha)<\ell_j(\alpha)\text{ for every }j\ne i\}.
\]
Ties are excluded. Write
\[
B_A(\ell)=\sum_{i=1}^k\bigl(1+\dim\operatorname{Aff}_{\R}C_i(\ell)\bigr),
\]
using the conventional empty-cell contribution zero, equivalently $\dim\operatorname{Aff}_{\R}\varnothing=-1$. The supplied version of Theorem 7 gives the tropical lower bound $\dim(kC)\geq B_A(\ell)$. The reconstructed universal question is whether, for every admissible exponent set $A$ and every $k\geq1$,
\[
\dim(kC)=\max_{\ell=(\ell_1,\ldots,\ell_k)}B_A(\ell).
\]
The dimension on the left is the algebraic dimension of the affine secant cone, not the projective secant dimension.

The cleaned statement instead calls $C$ simply the closure of the affine image of $f$. Identifying that closure with a cone needs a homogeneity convention, such as exponents in a hyperplane of nonzero height; merely saying ``an affine hyperplane'' does not state that condition. The original cone convention and the cleaned affine-image convention must therefore be checked rather than silently identified. The empty-cell convention and field hypotheses also need source confirmation.

The registry labels the formulation reconstructed and unverified. Its verification note says the report supplies the monomial setup, selected subsets and explicit bound, but that the equality question is a conservative rendering of a request to optimize the bound. Its triage dated 2026-08-21 labels the question open and reports exactness in numerous low-dimensional homogeneous and toric examples without claiming universal equality.
\subsection{Short English statement}
For a monomial cone, optimize the stated tropical lower bound over all tuples of affine-linear forms and ask whether the result always equals its secant-cone dimension. The universal equality is a reconstructed interpretation, and cone, field and empty-cell conventions need confirmation.
\subsection{Sources}
\begin{itemize}
\item Local UnsolvedMath registry, permanent ID 30000891, source code catalogue entry; supplied complete record, including the reconstructed-unverified qualification and reference to Theorem 7.
\item Tropical Geometry (2008), Oberwolfach Reports, DOI 10.4171/owr/2007/57, pp. 3303--3368: \url{https://doi.org/10.4171/owr/2007/57}. Proposed-year metadata gives 2007.
\item Tropical approach to secant dimensions: \url{https://arxiv.org/abs/math/0605345}.
\item Low-dimensional homogeneous secants: \url{https://doi.org/10.1515/advgeom.2010.001}.
\end{itemize}
\end{document}
